\begin{titlepage}
  \centering
  \vspace*{2cm}
  {\LARGE\bfseries Дискретная информационная\\[0.4em] термодинамика Вселенной\par}
  \vspace{1.2cm}
  {\large Монография о решёточной жидкости на планковском масштабе:\\
  от локального закона эволюции к наблюдаемой физике\par}
  \vfill
  {\large \today\par}
\end{titlepage}

\cleardoublepage
\thispagestyle{empty}
\vspace*{\fill}
\begin{center}
  \itshape
  Моей учительнице физики ---\\[0.6em]
  ООШ №21, г.~Киселёвск,\\[0.3em]
  выпуск 2009.\\[1.4em]
  \normalfont\small
  На лабораторных она требовала проверять по размерностям\\
  каждую формулу. С математикой у меня всегда было нелегко\\
  (СДВГ, потерянный знак, уехавший коэффициент) ---\\
  а размерность чаще всего ловит ошибку раньше,\\
  чем она разрастается. Это не занудство, а спасение.
\end{center}
\vspace*{\fill}
\cleardoublepage

\chapter*{Предисловие}
\addcontentsline{toc}{chapter}{Предисловие}

Спуск одной лестницей.
Сначала называются требования физики.
Затем они же записываются как макромир, микромир, квантовая механика и теория поля.
Там, где непрерывное поле в точке перестаёт их выполнять, те же требования становятся планковским миром --- моделью.
Дальше реестр условий, геометрия, закон эволюции и возвращение к макромиру идут из этого же списка.

Изложение --- как в теоретическом учебнике: определения, теоремы и доказательства
связаны перекрёстными ссылками.
Где уместно, за выкладками следит размерностной анализ --- школьная привычка
с физической лаборатории, без которой эта монография выглядела бы иначе.
Краткие пояснения по обратимым дискретным схемам и представлению Маделунга ---
в приложении~\ref{app:background}.

Читатель, знакомый с квантовой механикой, термодинамикой и специальной теорией относительности,
может идти по тексту последовательно: требования физики, перенос вниз, теоремы, восстановление известной физики.

% Печать notation-glossary (записи: glossary/notation-entries.tex).
\chapter*{Условные обозначения}
\addcontentsline{toc}{chapter}{Условные обозначения}
\label{ch:notation}

% Все записи type=notation, порядок по полю sort (не noidx: иначе пусто до «следующего» прогона).
\printunsrtglossary[type=notation,style=notationlong,sort=standard,nonumberlist]


Список основных обозначений --- выше (решётка $a$, $\Delta t$; длина $\lambda$; время $\tau$, $t_{\mathrm{SI}}$; кольцо $N_{\mathrm{ring}}$).

\tableofcontents
